% Figaro: Self-Enforcing Multi-Party Coordination via Asymmetric Bonding
% Prepared for Financial Cryptography and Data Security (FC)
% April 2026

\documentclass[11pt,a4paper]{article}

% ── Packages ────────────────────────────────────────────────────────
\usepackage[utf8]{inputenc}
\usepackage[T1]{fontenc}
\usepackage{amsmath,amssymb,amsthm}
\usepackage{mathtools}
\usepackage{algorithm}
\usepackage{algpseudocode}
\usepackage{booktabs}
\usepackage{graphicx}
\usepackage{hyperref}
\usepackage{cleveref}
\usepackage{enumitem}
\usepackage{xcolor}
\usepackage{listings}
\usepackage[margin=1in]{geometry}
\usepackage{natbib}
\bibliographystyle{plainnat}

% ── Theorem environments ────────────────────────────────────────────
\newtheorem{theorem}{Theorem}[section]
\newtheorem{lemma}[theorem]{Lemma}
\newtheorem{corollary}[theorem]{Corollary}
\newtheorem{proposition}[theorem]{Proposition}
\theoremstyle{definition}
\newtheorem{definition}[theorem]{Definition}
\newtheorem{example}[theorem]{Example}
\theoremstyle{remark}
\newtheorem{remark}[theorem]{Remark}

% ── Listings config ─────────────────────────────────────────────────
\lstset{
  basicstyle=\ttfamily\small,
  keywordstyle=\bfseries\color{blue!70!black},
  commentstyle=\itshape\color{gray},
  frame=single,
  breaklines=true,
  captionpos=b,
}

% ── Macros ──────────────────────────────────────────────────────────
\newcommand{\figaro}{\textsc{Figaro}}
\newcommand{\pay}{P}           % payment
\newcommand{\cumval}{G}        % cumulative value
\newcommand{\sbond}{C_s}       % seller bond
\newcommand{\bbond}{C_b}       % buyer bond
\newcommand{\spay}{\pi_s}      % seller payout
\newcommand{\bpay}{\pi_b}      % buyer payout
\newcommand{\ub}{u_B}          % buyer utility
\newcommand{\us}{u_S}          % seller utility
\DeclareMathOperator{\Coop}{C}
\DeclareMathOperator{\Def}{D}

% ────────────────────────────────────────────────────────────────────

\title{%
  \figaro: Self-Enforcing Multi-Party Coordination\\
  via Asymmetric Bonding%
}

\author{
  Alessandro Daliana\thanks{Corresponding author. \texttt{adaliana@figaro.org}}
}

\date{April 2026}

\begin{document}

\maketitle

% ════════════════════════════════════════════════════════════════════
\begin{abstract}
We present \figaro, a smart-contract protocol that enables self-enforcing
agreements between mutually distrusting parties without arbitrators,
timeouts, or governance mechanisms. The core primitive is
\emph{asymmetric bonding}: each party locks collateral equal to twice the
transaction value; only the buyer can trigger resolution. Cooperation is
thereby the strictly dominant strategy in a one-shot game, producing a
Nash equilibrium that holds without repeated interaction.

The mechanism scales from two-party exchange to $N$-party process trees
through \emph{progressive collateralization}: each seller bonds against
the cumulative value of all upstream work, not merely their own payment.
We prove that this preserves the Nash equilibrium at every chain
position, with the risk-to-reward ratio increasing geometrically for
downstream participants.

All orders within a process resolve \emph{atomically}---all or
nothing---creating endogenous coordination pressure among sellers
analogous to group lending in microfinance. This yields a three-layer
enforcement architecture: (1)~economic self-interest via bonding,
(2)~inter-seller coordination pressure via atomic resolution, and
(3)~legal deterrence backed by immutable on-chain evidence.

We describe the mechanism formally, prove its game-theoretic properties,
present its implementation as a Solidity smart contract verified
by TLA\textsuperscript{+} model checking (7~invariants, 2M+ states) and property-based fuzzing, and
discuss implications for organizational theory---specifically, the
Coasean thesis that firms exist to minimize transaction costs, arguing
that sufficiently low enforcement cost shifts coordination toward
transaction-scoped institutions assembled directly from bonded
counterparties.

\end{abstract}

\medskip
\noindent\textbf{Keywords:}
mechanism design, Nash equilibrium, smart contracts, collateral bonding,
multi-party coordination, process trees, post-firm economy

% ════════════════════════════════════════════════════════════════════
\section{Introduction}
\label{sec:intro}

\subsection{The Problem}

Economic coordination between strangers requires trust infrastructure.
When two parties who have never met wish to exchange value, they face the
classic commitment problem: either party may defect after the other has
performed. The historical solution is \emph{institutional
intermediation}---firms, courts, escrow services, platform
operators---each of which introduces costs, delays, jurisdictional
dependencies, and power asymmetries.

The problem intensifies in multi-party settings. A food delivery involves
a cook, a kitchen operator, an ingredient sourcer, and a courier; a
procurement process involves engineers, manufacturers, inspectors, and
logistics providers. Traditional coordination architectures address this
complexity by aggregating participants into \emph{firms}---legal entities
that internalize transaction costs at the expense of individual
sovereignty. Coase~\citep{coase1937} formalized this observation:
firms exist because the cost of discovering prices, negotiating terms,
and enforcing contracts in open markets exceeds the cost of hierarchical
management.

Blockchain-based smart contracts have produced a rich literature on
trustless exchange, but most protocols address financial
operations---swaps, lending, derivatives---rather than general
coordination. Those that target coordination
(e.g.,~\citealt{kleros2019,aragon2017}) typically reintroduce the
intermediary as an on-chain dispute resolution mechanism: a jury, a
governance body, or a timeout that unilaterally releases funds.

\subsection{Our Contribution}

We present \figaro{}, which achieves self-enforcing coordination through
a single mechanism: \emph{asymmetric bonding}. The protocol's
contributions are:

\begin{enumerate}[label=(\roman*)]
  \item \textbf{A minimal settlement primitive} with two external
    functions, three mappings, no owner, no fee, and no escape hatches
    from the bonded state. We show that this minimality is not a
    simplification but a \emph{security requirement}: every additional
    exit path weakens the Nash equilibrium.

  \item \textbf{Progressive collateralization} for $N$-party process
    trees. We prove (\Cref{thm:n-party-nash}) that cooperation remains
    strictly dominant at every position in an arbitrary-depth process
    tree, with the seller's risk-to-reward ratio increasing with chain
    depth.

  \item \textbf{Atomic resolution} as a coordination forcing function.
    We formalize the analogy between all-or-nothing process settlement
    and group lending in microfinance, showing that atomic resolution
    creates endogenous peer pressure among sellers without explicit
    communication or governance.

  \item \textbf{Self-enforcing DAG topology} without on-chain
    validation. We prove (\Cref{thm:self-enforcing-dag}) that
    misreporting cumulative value is a strictly dominated strategy,
    making on-chain parent-relationship validation unnecessary.

  \item \textbf{Formal verification.} The implementation is verified
    through TLA\textsuperscript{+} model checking (7~invariants,
    2M+ states exhaustively explored), property-based fuzzing (7~Echidna
    invariants, 50{,}000 call sequences), and 180~Foundry unit tests.

  \item \textbf{Implications for organizational theory.} We argue that
    asymmetric bonding prices the cost of trust at exactly $2\times$ the
    transaction value, potentially replacing the standing firm with
    transaction-scoped institutions assembled directly around each
    process.
\end{enumerate}

\subsection{Paper Organization}

\Cref{sec:related} surveys related work. \Cref{sec:mechanism} presents
the mechanism formally. \Cref{sec:game-theory} proves the game-theoretic
properties. \Cref{sec:process-trees} extends the model to $N$-party
process trees and DAGs. \Cref{sec:implementation} describes the smart
contract implementation and verification. \Cref{sec:security} analyzes
attack vectors. \Cref{sec:post-firm} discusses implications for
organizational economics. \Cref{sec:conclusion} concludes with open
questions and future work.

% ════════════════════════════════════════════════════════════════════
\section{Related Work}
\label{sec:related}

\paragraph{Mutual-destruction escrows.}
The direct ancestor of \figaro{} is the Safe Remote Purchase
contract~\citep{buterin2016}, included in the Solidity documentation as
a reference pattern: buyer and seller each lock $2\times$ the item
value, and only mutual agreement releases the funds. This two-party,
single-order pattern is the seed from which \figaro{} grew.
BitHalo~\citep{bithalo2014} independently explored mutual-destruction
deposits for peer-to-peer Bitcoin trade.

The central problem was generalizing SRP beyond two parties. SRP uses
symmetric $2\times$ deposits with mutual agreement to release---both
parties must cooperate. This works for a single exchange but cannot
compose: in a twelve-seller delivery workflow, there is no natural
``mutual agreement'' among all participants. \emph{Asymmetric bonding}
(\Cref{def:asymmetric-bonding}) solves this. By scaling each seller's
bond with cumulative upstream value rather than fixing it per-trade,
the mechanism creates a mesh of independently secured edges. Each edge
carries its own Nash equilibrium (\Cref{thm:two-party-nash}), and the
mesh decomposes as finely as gas costs permit.

\emph{Buyer dominance} then takes advantage of this mesh structure to
make the standing firm non-compulsory. Because a single party resolves
the entire process tree atomically, coordination pressure propagates
through the mesh without any management layer
(\Cref{prop:coordination-pressure}). What was previously internalized
inside a firm---scheduling, quality control, supplier
management---is now enforced by the bonding game itself. The
organizational unit becomes process-local and temporary rather than a
standing managerial container.

Extending to $N$ parties followed naturally once buyer dominance was
established---the cumulative value accumulator and atomic resolution
are direct consequences. The difficult work was \emph{reduction}:
removing timeouts, owner functions, governance mechanisms, internal
ledgers, withdrawal patterns, position NFTs, and every other feature
that development pressure---from collaborators, AI agents, and
inherited DeFi/web2 mental models---insisted on adding. The result is
a minimal kernel: two external functions, three mappings, no owner, no
fee, no escape hatches. Each removed feature was not merely unnecessary
but actively harmful---violating one of the six properties derived in
\Cref{sec:game-theory}.

\paragraph{State channels and payment networks.}
The Lightning Network~\citep{poon2016} and Raiden~\citep{raiden2017} use
time-locked deposits for off-chain payment. These protocols optimize for
high-frequency bilateral payments and rely on timeouts for dispute
resolution---precisely the escape hatch that \figaro{} eliminates. More
critically, state channels do not address multi-party \emph{coordination}
(delivery, fulfillment, service composition); they address multi-hop
\emph{payment routing}.

\paragraph{On-chain dispute resolution.}
Kleros~\citep{kleros2019} provides decentralized arbitration via
Schelling-point juror selection. Aragon Court~\citep{aragon2017} offers
similar functionality within DAO governance. Both reintroduce a
\emph{discretionary human decision} into the resolution path, breaking
the self-enforcing property. In \figaro{}, disputes are resolved by
economic pressure (Layer~1), peer coordination (Layer~2), and---only in
extremis---off-chain legal proceedings backed by immutable on-chain
evidence (Layer~3). No on-chain governance mechanism exists.

\paragraph{Optimistic mechanisms.}
Optimistic rollups~\citep{optimism2021} and optimistic
oracles~\citep{uma2020} use a challenge-response pattern: an assertion
stands unless challenged within a timeout window. This requires
liveness assumptions (the challenger must be online) and
timeout calibration (too short enables fraud; too long delays
finality). \figaro{} has no timeout from the bonded state---capital
remains locked indefinitely until the buyer resolves, making liveness
irrelevant to security.

\paragraph{Mechanism design.}
The revelation principle~\citep{myerson1981} guarantees that any
achievable outcome can be implemented via a direct mechanism. Figaro's
insight is that the ``direct mechanism'' for multi-party coordination
need not involve message passing, reporting, or arbitration at all:
locked capital \emph{is} the mechanism. The closest theoretical
antecedent is the concept of \emph{deposit
games}~\citep{asgaonkar2019}, which analyze two-party security deposits.
We extend deposit game analysis to $N$-party DAGs with progressive
collateralization.

\paragraph{Transaction cost economics.}
Coase~\citep{coase1937} argued that firms exist to minimize the
transaction costs of market exchange. Williamson~\citep{williamson1985}
refined this into Transaction Cost Economics (TCE), identifying asset
specificity, uncertainty, and frequency as drivers of hierarchical
governance. We argue (\Cref{sec:post-firm}) that asymmetric bonding
provides a \emph{fixed-cost} alternative to these institutional
solutions, potentially altering the make-or-buy boundary that defines
the firm.

% ════════════════════════════════════════════════════════════════════
\section{The Mechanism}
\label{sec:mechanism}

\subsection{Notation and Setup}

Consider two parties, a \emph{buyer} $B$ and a \emph{seller} $S$, who
wish to exchange a payment $\pay > 0$ (denominated in an ERC-20 token)
for goods or services. Both parties hold sufficient token balances.

\begin{definition}[Asymmetric Bonding]
\label{def:asymmetric-bonding}
An \emph{asymmetric bond} for a transaction with payment $\pay$ and
cumulative value $\cumval \geq \pay$ is a pair of deposits:
\begin{align}
  \bbond &= 2\pay \quad (\text{buyer bond}), \label{eq:buyer-bond} \\
  \sbond &= 2\cumval \quad (\text{seller bond}). \label{eq:seller-bond}
\end{align}
For a root transaction (no upstream value), $\cumval = \pay$, so both
bonds equal $2\pay$. For sub-orders in a process tree,
$\cumval > \pay$, creating an asymmetry where the seller's bond exceeds
the buyer's.
\end{definition}

\begin{definition}[Commitment]
\label{def:commitment}
A \emph{commitment} is a tuple $\langle B, S, \tau, \pay, h, \sigma,
\delta \rangle$ where $B$ is the buyer address, $S$ the seller address,
$\tau$ the token contract, $\pay$ the payment, $h$ a content hash of
the off-chain agreement, $\sigma$ a cryptographic salt, and $\delta$ a
signature expiry deadline. Both parties sign the commitment off-chain
using EIP-712 typed structured data~\citep{eip712}.
\end{definition}

\begin{definition}[Process]
\label{def:process}
A \emph{process} $\Pi$ is a tuple $\langle \text{id}, B, \tau, \cumval,
n, \mathcal{O} \rangle$ where $\text{id}$ is a deterministic identifier
derived from the root commitment, $B$ is the root buyer (constant across
all orders), $\tau$ is the denomination token (constant), $\cumval$ is a
\emph{monotonic accumulator} tracking total committed value, $n$ is the
count of active orders, and $\mathcal{O}$ is the set of active order
identifiers.
\end{definition}

\subsection{Protocol Operations}

The protocol exposes exactly two state-changing operations.

\paragraph{Operation 1: \texttt{commit}.}
Given a commitment $c$ signed by both $B$ and $S$:
\begin{enumerate}
  \item Verify both EIP-712 signatures.
  \item If $c.\text{processId} = 0$, derive a new process identifier from
    the typed-data digest and initialize the process with
    $\cumval \gets \pay$ and $n \gets 1$.
  \item Otherwise verify that $c.\text{processId}$ identifies an
    existing process, that signer $B$ matches the process root buyer
    (\emph{buyer dominance}), and that
    $c.\text{expectedCumulativeValue} = \Pi.\cumval + \pay$.
  \item Pull $\bbond = 2\pay$ from $B$ and
    $\sbond = 2c.\text{expectedCumulativeValue}$ from $S$.
  \item Create or extend the process accordingly, update the monotonic
    accumulator, and emit an evidence event that fully records the
    commitment.
\end{enumerate}

\paragraph{Operation 2: \texttt{resolveProcess}.}
Given a process identifier and the complete list of active commitments:
\begin{enumerate}
  \item Verify $\texttt{msg.sender} = \Pi.B$ (buyer dominance).
  \item Verify that the list length equals $\Pi.n$ (atomic resolution:
    \emph{all} orders must be included).
  \item For each order $o_i$ with payment $\pay_i$ and cumulative snapshot
    $\cumval_i$:
    \begin{align}
      \pi_{s,i} &= 2\cumval_i + \pay_i, \label{eq:seller-payout} \\
      \pi_{b,i} &= \pay_i. \label{eq:buyer-payout}
    \end{align}
  \item Transfer $\pi_{s,i}$ to each seller, $\pi_{b,i}$ to the buyer.
  \item Mark all orders as resolved.
\end{enumerate}

\begin{remark}[Conservation Law]
\label{rem:conservation}
For each order, the total distributed equals total locked:
\begin{equation}
  \pi_{s,i} + \pi_{b,i} = (2\cumval_i + \pay_i) + \pay_i = 2\cumval_i + 2\pay_i = C_{s,i} + C_{b,i}.
\end{equation}
The contract holds zero balance after resolution.
\end{remark}

\begin{remark}[No Escape Hatches]
\label{rem:no-escape}
There is no timeout, no cancel-after-commitment, no admin override, no
governance vote, and no unilateral exit from the bonded state. Once both
parties have committed, the only path to fund recovery is buyer-initiated
resolution. This is not a temporary simplification; it is a
\emph{security requirement.} \Cref{thm:escape-hatch-weakness} proves
that any additional exit path weakens the Nash equilibrium.
\end{remark}

\subsection{Net Payoffs}

\begin{definition}[Generalized Payout Functions]
\label{def:payout-functions}
For a bond multiplier $k \geq 1$, an order with payment $\pay > 0$
and cumulative value $\cumval \geq \pay$, define the
\emph{settlement payout functions}:
\begin{align}
  \spay(k, \cumval, \pay) &= k \cdot \cumval + \pay,
    \label{eq:general-seller-payout} \\
  \bpay(k, \pay) &= (k - 1) \cdot \pay.
    \label{eq:general-buyer-payout}
\end{align}
These satisfy the conservation law
$\spay + \bpay = k \cdot \cumval + k \cdot \pay = \sbond + \bbond$
for $\sbond = k \cdot \cumval$ and $\bbond = k \cdot \pay$. The
\emph{net payoff} is the payout minus the deposited bond:
\begin{align}
  \us^{\text{net}}(k) &= \spay - \sbond = \pay,
    \label{eq:seller-net} \\
  \ub^{\text{net}}(k) &= \bpay - \bbond = -\pay.
    \label{eq:buyer-net}
\end{align}
\end{definition}

\begin{remark}
For the protocol's $k = 2$: $\spay = 2\cumval + \pay$,
$\bpay = \pay$, $\sbond = 2\cumval$, $\bbond = 2\pay$. The
net payoffs ($+\pay$ for the seller, $-\pay$ for the buyer) are
independent of $k$---the multiplier affects only the capital at risk,
not the final value transfer.
\end{remark}

% ════════════════════════════════════════════════════════════════════
\section{Game-Theoretic Analysis}
\label{sec:game-theory}

\subsection{Two-Party Nash Equilibrium}

\begin{definition}[The Bonding Game]
\label{def:bonding-game}
The \emph{bonding game} $\Gamma = (\{B, S\}, \{A_B, A_S\}, \{\ub, \us\})$
is a one-shot, simultaneous-move, complete-information game:
\begin{itemize}
  \item \textbf{Players}: buyer $B$, seller $S$.
  \item \textbf{Action sets}: $A_B = A_S = \{\Coop, \Def\}$.
  \item \textbf{Interpretation}: $\Coop$ for $S$ means delivering the
    agreed goods/services; $\Coop$ for $B$ means calling
    \texttt{resolveProcess}. $\Def$ means failing to perform.
  \item \textbf{Payoff functions} (net of bond deposits), for a root
    order with payment $\pay > 0$ and $\cumval = \pay$:
\end{itemize}
\begin{equation}\label{eq:utility-buyer}
  \ub(a_B, a_S) =
  \begin{cases}
    -\pay & \text{if } a_B = \Coop \text{ and } a_S = \Coop, \\
    -2\pay & \text{otherwise}.
  \end{cases}
\end{equation}
\begin{equation}\label{eq:utility-seller}
  \us(a_B, a_S) =
  \begin{cases}
    +\pay & \text{if } a_B = \Coop \text{ and } a_S = \Coop, \\
    -2\cumval & \text{otherwise}.
  \end{cases}
\end{equation}
The ``otherwise'' payoff is identical across all non-cooperative
outcomes because any defection prevents resolution: bonds remain locked
indefinitely, yielding a net loss equal to the full deposit.
\end{definition}

\begin{table}[ht]
\centering
\caption{Payoff matrix for the bonding game $\Gamma$.
Entry format: $(\ub, \us)$. Parameters: $\pay > 0$,
$\cumval \geq \pay > 0$.}
\label{tab:payoff-matrix}
\begin{tabular}{lcc}
\toprule
 & $a_B = \Coop$ & $a_B = \Def$ \\
\midrule
$a_S = \Coop$ & $(-\pay,\; +\pay)$ & $(-2\pay,\; -2\cumval)$ \\
$a_S = \Def$  & $(-2\pay,\; -2\cumval)$ & $(-2\pay,\; -2\cumval)$ \\
\bottomrule
\end{tabular}
\end{table}

\begin{definition}[Dominance]
\label{def:dominance}
Action $a_i$ \emph{weakly dominates} $a_i'$ for player $i$ if
$u_i(a_i, a_{-i}) \geq u_i(a_i', a_{-i})$ for all $a_{-i}$, with
strict inequality for at least one $a_{-i}$. If the inequality is
strict for \emph{all} $a_{-i}$, the dominance is \emph{strict}.
\end{definition}

\begin{theorem}[Two-Party Nash Equilibrium]
\label{thm:two-party-nash}
In the bonding game $\Gamma$ of \Cref{def:bonding-game}:
\begin{enumerate}[label=(\alph*)]
  \item $\Coop$ weakly dominates $\Def$ for both the buyer and
    the seller.
  \item $(\Coop, \Coop)$ is the unique Nash equilibrium.
  \item $(\Coop, \Coop)$ is the unique Pareto-optimal outcome.
\end{enumerate}
\end{theorem}

\begin{proof}
\textbf{Part (a): Weak dominance.}

\emph{Buyer.} We compare $\ub(\Coop, a_S)$ vs.\
$\ub(\Def, a_S)$ for each seller action:
\begin{align}
  a_S = \Coop: &\quad \ub(\Coop, \Coop) = -\pay > -2\pay = \ub(\Def, \Coop),
    \label{eq:buyer-strict} \\
  a_S = \Def: &\quad \ub(\Coop, \Def) = -2\pay = \ub(\Def, \Def).
    \label{eq:buyer-weak}
\end{align}
Inequality~\eqref{eq:buyer-strict} is strict (since $\pay > 0$);
\eqref{eq:buyer-weak} is an equality. By \Cref{def:dominance}, $\Coop$
weakly dominates $\Def$ for $B$.

\emph{Seller.} We compare $\us(a_B, \Coop)$ vs.\
$\us(a_B, \Def)$ for each buyer action:
\begin{align}
  a_B = \Coop: &\quad \us(\Coop, \Coop) = +\pay > -2\cumval = \us(\Coop, \Def),
    \label{eq:seller-strict} \\
  a_B = \Def: &\quad \us(\Def, \Coop) = -2\cumval = \us(\Def, \Def).
    \label{eq:seller-weak}
\end{align}
Inequality~\eqref{eq:seller-strict} is strict (since $\pay > 0$ and
$\cumval > 0$ imply $\pay + 2\cumval > 0$);
\eqref{eq:seller-weak} is an equality. $\Coop$ weakly dominates $\Def$
for $S$.

\textbf{Part (b): Unique Nash equilibrium.}
Since $\Coop$ weakly dominates $\Def$ for both players, neither can
profitably deviate from $(\Coop, \Coop)$. We verify no other profile
is a Nash equilibrium:
\begin{itemize}
  \item $(\Coop, \Def)$: Seller deviates to $\Coop$, gaining
    $\pay - (-2\cumval) = \pay + 2\cumval > 0$. Not a NE.
  \item $(\Def, \Coop)$: Buyer deviates to $\Coop$, gaining
    $-\pay - (-2\pay) = \pay > 0$. Not a NE.
  \item $(\Def, \Def)$: Buyer deviates to $\Coop$ with no change
    (both yield $-2\pay$). Seller deviates to $\Coop$ with no change
    (both yield $-2\cumval$). This profile satisfies the Nash
    condition (no \emph{profitable} deviation), but an iterated
    elimination of weakly dominated strategies removes $\Def$ for
    both players, leaving $(\Coop, \Coop)$ as the unique survivor.
\end{itemize}

\textbf{Part (c): Pareto optimality.}
The payoff vector $(-\pay, +\pay)$ Pareto-dominates $(-2\pay, -2\cumval)$
since $-\pay > -2\pay$ and $+\pay > -2\cumval$.
\end{proof}

\begin{remark}[On Weak vs.\ Strict Dominance]
\label{rem:dominance-note}
The dominance is weak, not strict, because when the counterparty
defects, both actions yield the same locked-capital payoff. This is
intrinsic to the mechanism: once one party defects, the outcome is
fully determined regardless of the other's choice. The critical
property is that $(\Def, \Def)$ does not survive iterated elimination
of weakly dominated strategies---the standard refinement for
mechanism design~\citep{myerson1981}.
\end{remark}

\subsection{Sufficiency of the $2\times$ Multiplier}

We now characterize the multiplier $k$ in the generalized bonding game
where $\sbond = k \cdot \cumval$, $\bbond = k \cdot \pay$, and
payouts are given by \Cref{def:payout-functions}.

\begin{lemma}[Buyer Indifference at $k = 1$]
\label{lem:k1-buyer}
For $k = 1$, the buyer is indifferent between cooperating and defecting
when the seller cooperates:
$\ub(\Coop, \Coop) = \ub(\Def, \Coop) = -\pay$.
\end{lemma}
\begin{proof}
With $k = 1$: $\bbond = \pay$ and $\bpay = (k-1)\pay = 0$ (from
\eqref{eq:general-buyer-payout}). Under our accounting convention
(\Cref{def:payout-functions}), the buyer's net payoff under cooperation
is $\ub(\Coop, \Coop) = \bpay - \bbond = 0 - \pay = -\pay$. Under
defection, the buyer's bond remains locked:
$\ub(\Def, \Coop) = -\bbond = -\pay$. Since
$\ub(\Coop, \Coop) = \ub(\Def, \Coop) = -\pay$, the buyer has no
strictly profitable deviation from $\Def$---the mechanism provides
zero marginal incentive to resolve.
\end{proof}

\begin{theorem}[Minimal Viable Bond Multiplier]
\label{thm:2x-sufficient}
Let $k \in \mathbb{Z}_{\geq 1}$ be the bond multiplier. Then:
\begin{enumerate}[label=(\alph*)]
  \item For $k = 1$, $\Coop$ does not weakly dominate $\Def$ for the
    buyer: the buyer is indifferent when the seller cooperates.
  \item For $k = 2$, $\Coop$ weakly dominates $\Def$ for both players
    (\Cref{thm:two-party-nash}).
  \item For $k > 2$, the equilibrium properties are identical to $k = 2$,
    but capital requirements increase by $(k - 2)(\cumval + \pay)$ per
    order with no improvement in dominance relations.
\end{enumerate}
Thus $k = 2$ is the minimal integer multiplier that yields weak
dominance of $\Coop$ for both players, and is Pareto-optimal in
capital efficiency among all valid multipliers.
\end{theorem}

\begin{proof}
\textbf{Part (a).} See \Cref{lem:k1-buyer}. With $k = 1$, the buyer
recovers $\bpay = 0$ upon resolution, having deposited $\bbond = \pay$.
The net monetary payoff from cooperating is $-\pay$; the net from
defecting is also $-\pay$ (bond locked). The buyer cannot be motivated
by the mechanism alone to resolve.

\textbf{Part (b).} With $k = 2$: $\bpay = \pay$, $\bbond = 2\pay$.
Cooperation yields $\ub = -\pay$; defection yields $\ub = -2\pay$.
Since $-\pay > -2\pay$ (as $\pay > 0$), the buyer strictly prefers
cooperation when the seller cooperates. By
\Cref{thm:two-party-nash}(a), this establishes weak dominance.

\textbf{Part (c).} For $k \geq 2$:
$\ub(\Coop, \Coop) = -\pay$ and $\ub(\Def, \Coop) = -k\pay$.
Since $-\pay > -k\pay$ for all $k \geq 2$, the dominance relation
holds. The surplus capital locked, $(k - 2)\pay$ for the buyer and
$(k - 2)\cumval$ for the seller, earns zero return and is pure
deadweight cost.
\end{proof}

\subsection{Escape-Hatch Weakness}

\begin{theorem}[Escape-Hatch Weakness]
\label{thm:escape-hatch-weakness}
Let $\Gamma_E$ be the bonding game augmented with a unilateral exit
action $E$ available to player $i$, returning fraction $\alpha \in (0, 1]$
of player $i$'s bond. Then either:
\begin{enumerate}[label=(\alph*)]
  \item $\Coop$ no longer weakly dominates $\Def$ for at least one
    player (the Nash equilibrium of \Cref{thm:two-party-nash}
    may not survive), or
  \item $E$ requires a decision by a third party $J \notin \{B, S\}$
    whose incentives are not constrained by the bond structure.
\end{enumerate}
\end{theorem}

\begin{proof}
We analyze three exhaustive sub-cases of unilateral exit.

\emph{Case 1: Timeout (automatic exit after duration $T$).}
The action space becomes $A_B = \{\Coop, \Def, \text{Wait}\}$ where
$\text{Wait}$ means the buyer does not resolve and collects
$\alpha \cdot \bbond$ after time $T$. The buyer's payoff under
$\text{Wait}$ against a cooperating seller is:
\begin{equation}
  \ub(\text{Wait}, \Coop) = -\bbond + \alpha \cdot \bbond
    = -2\pay(1 - \alpha).
  \label{eq:timeout-payoff}
\end{equation}
Compare with cooperation: $\ub(\Coop, \Coop) = -\pay$. The buyer
prefers $\Coop$ to $\text{Wait}$ iff:
\begin{equation}
  -\pay > -2\pay(1 - \alpha)
  \quad\Longleftrightarrow\quad
  2(1 - \alpha) > 1
  \quad\Longleftrightarrow\quad
  \alpha < \tfrac{1}{2}.
  \label{eq:timeout-threshold}
\end{equation}
For $\alpha \geq \frac{1}{2}$, $\text{Wait}$ yields
$\ub \geq -\pay = \ub(\Coop, \Coop)$, so $\Coop$ no longer
(weakly) dominates $\text{Wait}$. The buyer may rationally wait out the
timeout instead of resolving. Even for $\alpha < \frac{1}{2}$, the
defection cost drops from $2\pay$ to $2\pay(1 - \alpha)$, weakening
the deterrent. The seller, anticipating the timeout, faces a modified
game where the buyer's commitment to resolve is no longer credible---the
timeout provides an alternative path to partial capital recovery,
undermining the "no escape" property that sustains the equilibrium.

\emph{Case 2: Arbitrator override.}
A third party $J$ can invoke $E$ on behalf of either player.
Formally, this extends the game to three players $\{B, S, J\}$.
The mechanism's self-enforcing property requires that the payoff
structure \emph{alone} determines the outcome. Introducing $J$
makes the outcome contingent on $J$'s action, and $J$ is not bonded:
$J$ deposits nothing, risks nothing, and has no mechanism-constrained
incentive to decide correctly. Whether $J$ is an individual
(arbitrator), an algorithm (oracle), or a collective (jury), the
resolution path now depends on an agent outside the bond structure.
This is condition (b).

\emph{Case 3: Governance vote.}
A special case of Case~2 where $J$ is replaced by a set of voters
$\{J_1, \ldots, J_m\}$ who collectively decide on $E$. This inherits
all problems of Case~2, plus: (i)~coordination failures among voters,
(ii)~plutocratic capture if voting power is token-weighted,
(iii)~voter apathy (rational ignorance). Each voter $J_k$ is unbonded.

In all three cases, either $\Coop$ ceases to weakly dominate all
alternatives (Case~1), or resolution depends on an unbonded external
agent (Cases 2--3). \qedhere
\end{proof}

\begin{corollary}\label{cor:no-timeout}
No timeout duration $T$ and no recovery fraction $\alpha > 0$ can be
added to the bonding game without weakening the buyer's incentive to
cooperate. In particular, full recovery ($\alpha = 1$) reduces the
game to a zero-cost option, making defection strictly dominant for
the buyer.
\end{corollary}

% ════════════════════════════════════════════════════════════════════
\section{$N$-Party Process Trees}
\label{sec:process-trees}

\subsection{Progressive Collateralization}

We now extend the two-party model to an arbitrary-length chain of
sellers.

\begin{definition}[Process Tree]
\label{def:process-tree}
A \emph{process tree} is a sequence of bonded orders
$\langle o_1, o_2, \ldots, o_n \rangle$ sharing a common buyer $B$ and
token $\tau$. Order $o_i$ has local payment $\pay_i > 0$ and cumulative
value:
\begin{equation}
  \cumval_i = \sum_{j=1}^{i} \pay_j.
  \label{eq:cumulative}
\end{equation}
The seller of $o_i$ posts bond $C_{s,i} = 2\cumval_i$; the buyer
posts bond $C_{b,i} = 2\pay_i$.
\end{definition}

\begin{definition}[$N$-Party Bonding Game]
\label{def:n-party-game}
The \emph{$N$-party bonding game}
$\Gamma_N = (\{B, S_1, \ldots, S_n\}, \{A_i\}, \{u_i\})$
extends $\Gamma$ to $n$ sellers. Each seller $S_i$ has action set
$\{\Coop, \Def\}$. The buyer has a single action conditioned on all
sellers: $B$ resolves iff all sellers cooperate. Payoffs:
\begin{equation}\label{eq:n-party-seller-utility}
  u_{S_i}(a_B, a_{S_1}, \ldots, a_{S_n}) =
  \begin{cases}
    +\pay_i & \text{if } a_B = \Coop \text{ and } a_{S_j} = \Coop
      \;\forall j, \\
    -2\cumval_i & \text{otherwise}.
  \end{cases}
\end{equation}
The buyer's payoff is:
\begin{equation}\label{eq:n-party-buyer-utility}
  \ub(a_B, a_{S_1}, \ldots, a_{S_n}) =
  \begin{cases}
    -\sum_{i=1}^{n} \pay_i & \text{if } a_B = \Coop \text{ and }
      a_{S_j} = \Coop \;\forall j, \\
    -\sum_{i=1}^{n} 2\pay_i & \text{otherwise}.
  \end{cases}
\end{equation}
\end{definition}

\begin{theorem}[$N$-Party Nash Equilibrium]
\label{thm:n-party-nash}
In $\Gamma_N$, cooperation is a weakly dominant strategy for every
seller $S_i$ at every position $i \in \{1, \ldots, n\}$, and
$(\Coop, \Coop, \ldots, \Coop)$ is the unique outcome surviving
iterated elimination of weakly dominated strategies.
\end{theorem}

\begin{proof}
Fix an arbitrary seller $S_i$. Consider the two cases for the remaining
players' joint action profile $a_{-i}$.

\emph{Case 1: All other sellers cooperate and buyer resolves.} Then:
\begin{align}
  u_{S_i}(\ldots, \Coop_i, \ldots) &= +\pay_i, \\
  u_{S_i}(\ldots, \Def_i, \ldots) &= -2\cumval_i.
\end{align}
Since $\pay_i > 0$ and $\cumval_i > 0$, we have
$\pay_i > -2\cumval_i$, i.e., $\pay_i + 2\cumval_i > 0$. The
cooperation payoff strictly exceeds the defection payoff.

\emph{Case 2: At least one other seller defects, or buyer defects.}
Then resolution does not occur regardless of $S_i$'s action:
\begin{equation}
  u_{S_i}(\ldots, \Coop_i, \ldots) = u_{S_i}(\ldots, \Def_i, \ldots)
    = -2\cumval_i.
\end{equation}
Payoffs are equal.

Combining both cases, $\Coop$ weakly dominates $\Def$ for $S_i$, with
strict inequality in Case~1. Since the choice of $i$ was arbitrary,
this holds for all $n$ sellers. The buyer's dominance follows from
\Cref{thm:two-party-nash}(a), generalized: the buyer's payoff under
cooperation ($-\sum \pay_i$) strictly exceeds the payoff under defection
($-2\sum \pay_i$) when all sellers cooperate.

Iterated elimination: in the first round, $\Def$ is eliminated for all
players (weakly dominated). The unique surviving profile is
$(\Coop, \ldots, \Coop)$.
\end{proof}

\begin{remark}[Strength of the $N$-Party Result]
Note that the seller's dominance condition,
$\pay_i + 2\cumval_i > 0$, holds with a margin that \emph{grows}
with chain depth (since $\cumval_i = \sum_{j \leq i} \pay_j$ is
monotonic). The cooperation-defection payoff gap for $S_i$ is:
\begin{equation}
  \Delta_i = \pay_i - (-2\cumval_i) = \pay_i + 2\cumval_i
    = \pay_i + 2\!\sum_{j=1}^{i} \pay_j \geq 3\pay_i.
\end{equation}
The gap is at least $3\pay_i$ (equality only at the root where
$\cumval_i = \pay_i$) and grows with every upstream payment. This
is the formal basis of the ``geometric coordination pressure'' in
\Cref{cor:geometric-pressure}.
\end{remark}

\begin{corollary}[Geometric Coordination Pressure]
\label{cor:geometric-pressure}
The risk-to-reward ratio for seller $S_i$ is:
\begin{equation}
  \rho_i = \frac{C_{s,i}}{\pay_i} = \frac{2\cumval_i}{\pay_i}
    = \frac{2\sum_{j=1}^{i} \pay_j}{\pay_i}.
\end{equation}
For equal payments ($\pay_j = p$ for all $j$), $\rho_i = 2i$, growing
linearly with chain depth. For typical value chains where early stages
contribute the majority of value ($\pay_1 \gg \pay_n$), $\rho_n$ grows
super-linearly.
\end{corollary}

\begin{example}
\label{ex:food-delivery}
Alice orders food (\$10) from Bob, who needs delivery (\$2) from
Charlie:
\begin{center}
\begin{tabular}{lccccc}
\toprule
Party & Role & $\pay_i$ & $\cumval_i$ & $C_{s,i}$ & $\rho_i$ \\
\midrule
Bob     & Seller 1 & \$10 & \$10 & \$20 & 2.0 \\
Charlie & Seller 2 & \$2  & \$12 & \$24 & 12.0 \\
\bottomrule
\end{tabular}
\end{center}
Charlie risks \$24 to earn \$2 ($\rho = 12$). Bob risks \$20 to earn
\$10 ($\rho = 2$). The deeper participant faces amplified consequences.
\end{example}

\subsection{Atomic Resolution and Seller Coordination}

\begin{definition}[Atomic Resolution]
\label{def:atomic-resolution}
Resolution is \emph{atomic}: the buyer must resolve all active orders
in a process simultaneously, or none. There is no partial resolution.
\end{definition}

\begin{proposition}[Endogenous Coordination Pressure]
\label{prop:coordination-pressure}
The $N$-party bonding game under atomic resolution induces a
\emph{weakest-link public goods game} among sellers, where each
seller's payout depends on universal cooperation.
\end{proposition}

\begin{proof}
Consider the sellers' subgame (fixing the buyer's strategy at: resolve
iff all sellers cooperate). Each $S_i$ has payoff
$u_{S_i}$ as in \eqref{eq:n-party-seller-utility}.

Define the \emph{externality} imposed by $S_k$'s defection on $S_i$
($i \neq k$):
\begin{equation}
  \varepsilon_{k \to i}
    = u_{S_i}(\text{all}\;\Coop) - u_{S_i}(S_k\;\text{defects})
    = \pay_i - (-2\cumval_i)
    = \pay_i + 2\cumval_i.
\end{equation}
This is strictly positive for all $i, k$ with $i \neq k$. Every
defection imposes a loss on every other seller. The total externality
of $S_k$'s defection is:
\begin{equation}
  \mathcal{E}_k = \sum_{i \neq k} \varepsilon_{k \to i}
    = \sum_{i \neq k} (\pay_i + 2\cumval_i).
\end{equation}

This structure satisfies the defining properties of a \emph{weakest-link
game}~\citep{hirshleifer1983}:
\begin{enumerate}
  \item The group outcome depends on the minimum individual effort
    (one defection collapses all payoffs to the non-cooperative level).
  \item The cooperative outcome $(\Coop, \ldots, \Coop)$ is the unique
    Pareto-optimal Nash equilibrium of the subgame.
  \item The game exhibits \emph{strategic complementarity}: $S_i$'s
    marginal return to cooperation is highest when all others cooperate.
\end{enumerate}

In a repeated-interaction setting, this externality structure creates
endogenous peer monitoring: $S_i$ has a direct economic interest of
magnitude $\pay_i + 2\cumval_i$ in ensuring $S_k$ cooperates. This
incentivizes information acquisition about co-sellers, pressure
application on underperformers, and reputation-based exclusion of known
defectors---all without explicit communication protocols or governance
mechanisms.
\end{proof}

\begin{remark}[Microfinance Analogy]
This structure mirrors group lending in microfinance
(e.g., Grameen Bank~\citep{grameen1999}), where default by one member of
a borrowing group causes the entire group to lose future credit access.
The microfinance literature documents substantially lower default rates
in group lending relative to individual lending, attributed to peer
monitoring and social collateral~\citep{ghatak1999}.
\end{remark}

\subsection{Self-Enforcing DAG Topology}

Process trees can express arbitrary directed acyclic graph (DAG)
topologies. A natural question: must the contract validate parent-child
relationships on-chain?

\begin{theorem}[Self-Enforcing DAG Honesty]
\label{thm:self-enforcing-dag}
In any game where seller $S_i$ chooses a reported cumulative value
$\cumval'_i \geq \cumval_i$ (with $\cumval_i$ being the true
value enforced by the contract's accumulator), truth-telling
($\cumval'_i = \cumval_i$) weakly dominates all alternatives.
The dominance is strict whenever the probability of non-resolution
is positive.
\end{theorem}

\begin{proof}
The contract's monotonic accumulator enforces
$\cumval'_i \geq \cumval_i \coloneqq \Pi.\cumval_{\text{prev}} + \pay_i$:
any report $\cumval'_i < \cumval_i$ causes a revert
(\texttt{CumulativeValueMismatch}), so we restrict attention to
$\cumval'_i \geq \cumval_i$.

Let $q \in [0, 1]$ be the probability that the buyer resolves the
process (a function of all sellers' performance, exogenous to $S_i$'s
reporting decision). Define $S_i$'s expected payoff as a function of
the report $\cumval'_i$:
\begin{equation}
  \mathbb{E}[u_{S_i}(\cumval'_i)]
    = q \cdot (\underbrace{2\cumval'_i + \pay_i}_{\text{payout}}
             - \underbrace{2\cumval'_i}_{\text{bond}})
    + (1 - q) \cdot (-2\cumval'_i)
    = q \cdot \pay_i - (1 - q) \cdot 2\cumval'_i.
  \label{eq:dag-expected-utility}
\end{equation}
The derivative with respect to $\cumval'_i$ is:
\begin{equation}
  \frac{\partial}{\partial \cumval'_i}
    \mathbb{E}[u_{S_i}] = -2(1 - q).
  \label{eq:dag-derivative}
\end{equation}
For $q < 1$ (non-zero probability of non-resolution), this derivative
is strictly negative: expected utility is strictly decreasing in the
reported cumulative value. The optimum is therefore at the smallest
feasible value, $\cumval'_i = \cumval_i$.

For $q = 1$ (certain resolution), the derivative is zero and the seller
is indifferent---overstating wastes capital but doesn't affect the net
payoff. Truth-telling remains weakly optimal.

Since $q < 1$ in any realistic setting (the seller cannot guarantee
other sellers' cooperation), truth-telling strictly dominates
overstating.
\end{proof}

% ════════════════════════════════════════════════════════════════════
\section{Implementation}
\label{sec:implementation}

\subsection{Smart Contract Architecture}

The protocol is implemented as a single Solidity smart contract,
\texttt{FigaroCore}, inheriting OpenZeppelin's EIP-712 and
ReentrancyGuard. The kernel exposes only two external functions:
\texttt{commit} and \texttt{resolveProcess}.

\paragraph{Storage layout.}
Three mappings comprise the kernel state:
\begin{itemize}
  \item \texttt{processes}: maps $\texttt{bytes32} \to
    \texttt{ProcessState}$ (root buyer, currency, cumulative value,
    active order count).
  \item \texttt{orderStatus}: maps $\texttt{bytes32} \to
    \texttt{uint8}$ (unknown, committed, resolved).
  \item \texttt{orderProcessId}: maps $\texttt{bytes32} \to
    \texttt{bytes32}$ (order membership in a process).
\end{itemize}
The kernel deliberately does not store a per-order struct. Order terms,
including payment and cumulative-value snapshot, are reconstructed from
the signed commitment and the emitted \texttt{OrderCommitted} event.
Bond amounts (\texttt{buyerBond = 2 $\times$ payment},
\texttt{sellerBond = 2 $\times$ cumulativeValue}) remain deterministic
functions of those values rather than separate storage variables.

\paragraph{Commitment protocol.}
Both parties sign an EIP-712 typed commitment off-chain. A single
on-chain \texttt{commit} call verifies both signatures and pulls both
bonds atomically. Root commitments create a new process; sub-order
commitments extend an existing one by carrying the inherited
\texttt{processId} and the next expected cumulative value. This
eliminates the accept-reject pattern of traditional escrow protocols and
prevents front-running of the counterparty's acceptance.

\paragraph{Identifier derivation.}
Process and order identifiers are content-addressed:
\begin{align}
  \text{processId} &= H_{\text{EIP-712}}(\text{commitment}) \quad \text{for root orders}, \\
  \text{orderHash} &= H(\text{processId} \| \text{structHash}).
\end{align}
For sub-orders, \texttt{processId} is inherited from the signed
commitment; the EIP-712 domain separator already binds the root
identifier to chain and verifying contract. This prevents collision,
replay, and cross-chain confusion without auto-incrementing counters.

\paragraph{Fee-on-transfer protection.}
The \texttt{\_pullExact} function checks
$\texttt{balanceAfter} - \texttt{balanceBefore} = \texttt{amount}$,
reverting if the token applies a transfer tax. This ensures the
conservation law (\Cref{rem:conservation}) cannot be broken by exotic
token behavior.

\subsection{Design Non-Decisions}

The following features are deliberately absent:

\begin{enumerate}
  \item \textbf{No owner or admin functions.} The contract is ownerless
    and has no \texttt{pause()}, \texttt{upgrade()}, or governance
    mechanism.
  \item \textbf{No protocol fee.} Settlement distributes the full bond
    amount to the parties. Any associated token layer is external to
    the kernel; in the current runtime, seller-only settlement
    emission serves as a coordination Schelling point rather than a
    fee path (\Cref{sec:conclusion}).
  \item \textbf{No timeout.} An active order remains active indefinitely.
    Per \Cref{thm:escape-hatch-weakness}, timeouts weaken the Nash
    equilibrium.
  \item \textbf{No partial resolution.} Per
    \Cref{prop:coordination-pressure}, partial resolution would destroy
    the seller coordination game.
  \item \textbf{No internal ledger.} Payouts are direct ERC-20 transfers,
    not balance increments. This eliminates withdrawal-pattern complexity
    and reentrancy surface.
\end{enumerate}

\subsection{Settlement-Anchored Emission}

The kernel itself has no fee and no token dependency. However, the
current runtime includes an associated coordination token (FIG) whose
emission is anchored to settlement events rather than calendar time.
This design is worth describing because the emission schedule is itself
a mechanism design artifact with properties that follow from the
protocol's structure.

Two emission paths coexist:

\paragraph{Direct path (on-chain).}
A permissionless, immutable contract (\texttt{FigEmission}) reads
\texttt{FigaroCore} order state to verify resolution, then mints FIG
to the seller. The rate is a discrete halving step function:
\begin{equation}
  r_{\text{direct}}(n) = r_0 \cdot 2^{-\lfloor n/H \rfloor},
  \label{eq:direct-emission}
\end{equation}
where $r_0 = 100$~FIG per settlement, $H = 10{,}000{,}000$~settlements
per halving epoch, and $n$ is the cumulative settlement count at the
time of the claim. A 600M~FIG hard cap terminates emission.

\paragraph{Batch path (SP1 prover).}
The batch path uses an Euler oscillation---a decaying cosine modulation
derived from Euler's identity---on top of the same decay envelope:
\begin{equation}
  r_{\text{batch}}(n) = r_0 \cdot 2^{-n/H} \cdot
    \bigl(1 + A \cdot \cos(2\pi n / S)\bigr),
  \label{eq:euler-emission}
\end{equation}
where $A = 0.5$ is the oscillation amplitude and
$S = 2{,}000{,}000$~settlements is the season length (one full
cosine cycle). The rate oscillates between $1.5 \times$ the decay
envelope (peaks) and $0.5 \times$ (troughs), creating ``emission
seasons'' whose calendar length is an emergent property of network
activity.

\begin{remark}[Zero-Integral Property]
The cosine term satisfies $\int_0^S \cos(2\pi n/S)\,dn = 0$. Over
each full season, the oscillation redistributes emission temporally
without changing the total---the Euler path and the direct path
converge in cumulative emission over long horizons.
\end{remark}

Three design choices warrant explanation. First, the decay is
settlement-count-based, not calendar-based. If adoption is slow, the
high early rate persists longer (the bootstrapping subsidy is still
needed); if adoption is fast, halving arrives quickly (network effects
have kicked in). Second, emission is seller-only: only the party
bearing the asymmetric capital commitment receives the reward. This
closes wash-trading vectors (a buyer-seller pair cannot profitably
cycle orders to extract FIG because each cycle requires locking $4\pay$
in real capital). Third, the batch oscillation creates temporal
coordination value: participants who batch settlements through the
SP1 prover during peak seasons earn up to $1.5\times$ the direct rate,
incentivizing the scalability path without penalizing the permissionless
direct fallback.

\subsection{Formal Verification}

\paragraph{TLA\textsuperscript{+} model checking.}
The contract is modeled in TLA\textsuperscript{+} with 3~actions
(\texttt{CommitRoot}, \texttt{CommitSub}, \texttt{ResolveProcess}) and
8~state variables. TLC verifies 7~invariants by exhaustive
enumeration over bounded state spaces (1~buyer, 2~sellers, payments
1--3, 2~concurrent processes, 2~sub-orders each; 2{,}095{,}093~states
generated, 1{,}521{,}909~distinct):
\begin{enumerate}[label=\textbf{I\arabic*}]
  \item \textbf{TypeOK}: type well-formedness of all state variables.
  \item \textbf{TokenConservation}: sum of all wallet balances plus
    contract balance equals total supply (conservation of value).
  \item \textbf{ContractSolvency}: contract balance $\geq 0$.
  \item \textbf{WalletNonNegative}: all participant balances $\geq 0$.
  \item \textbf{CumulativeIntegrity}: for every process,
    $\cumval$ equals the sum of all order payments.
  \item \textbf{ActiveCountCorrect}: stored active-order count
    matches the count of committed orders.
  \item \textbf{ResolutionAlwaysPossible}: for every active process,
    the contract holds sufficient funds to resolve all orders.
\end{enumerate}

The model abstracts EIP-712 signatures (assumed correct given valid
commitments), ERC-20 mechanics (modeled as integer balances), and
timing (deadlines are orthogonal to the bonding equilibrium).

\paragraph{Property-based fuzzing (Echidna).}
Seven property invariants are checked after every call sequence
(50{,}000~sequences of length~30, totaling $\sim$43{,}000 individual
calls):
\begin{enumerate}[label=\textbf{E\arabic*}]
  \item \textbf{Solvency}: token balance $\geq$ sum of outstanding bonds.
  \item \textbf{ActiveCountConsistent}: stored count matches actual.
  \item \textbf{CumulativeAccounting}: accumulator equals sum of payments.
  \item \textbf{StateMonotonicity}: orders never transition backwards.
  \item \textbf{BuyerDominance}: non-buyer resolve always reverts.
  \item \textbf{AtomicResolution}: incomplete order list always reverts.
  \item \textbf{NoSelfDeal}: $B = S$ always reverts.
\end{enumerate}
The Echidna fuzzer uses HEVM cheatcodes for EIP-712 signing with known
private keys, exercising the full commitment-to-resolution path.

\paragraph{Unit testing.}
236~Foundry unit tests cover the contract across 15~test files, including
EIP-712 signature verification, cumulative value enforcement, permit
integration, event emission, revert branch coverage, batch verifier
parity, FIG emission mechanics, and safety
hardening. An additional 166~tests cover the TypeScript SDK.

% ════════════════════════════════════════════════════════════════════
\section{Security Analysis}
\label{sec:security}

\subsection{Threat Model}

We assume rational adversaries who maximize economic payoff. We consider
four attack vectors.

\paragraph{Griefing (buyer refuses to resolve).}
The attacker's cost is $2\pay$ (locked bond) plus opportunity cost on
that capital plus permanent on-chain reputation damage. The attacker's
benefit is zero (no economic gain from non-resolution). For rational
agents, the attack is strictly dominated by cooperation. For irrational
agents (spite exceeds economic loss), Layer~3 provides deterrence: the
immutable on-chain record---including the order's commitment timestamp,
the buyer's address, and the fact that no resolution has been
called---serves as evidence in off-chain legal proceedings.

\paragraph{Sybil (fake orders to manipulate state).}
Each order requires real bond deposits:
$2(\cumval_i + \pay_i)$ total locked per order. At scale, Sybil attacks
are capital-intensive with no mechanism for extracting the invested
capital (there is no secondary market, no yield, and no governance
influence associated with locked bonds).

\paragraph{Front-running.}
Order identifiers are content-addressed from the signed commitment. An
attacker who observes a pending \texttt{commit} transaction cannot
create a conflicting order for the same process (different signatures
produce different hashes). No extractable MEV exists.

\paragraph{Cumulative value manipulation.}
Per \Cref{thm:self-enforcing-dag}, overstating cumulative value is
strictly dominated (greater capital at risk, same payoff). Understating
is prevented by the contract's monotonic accumulator check.

\subsection{Liveness}

\begin{theorem}[Liveness Under Rationality]
\label{thm:liveness}
Let $B$ be a buyer with time-discounted utility
$U_B = \sum_{t=0}^{\infty} \delta^t \cdot u_t$ where $\delta \in (0, 1)$
is the discount factor and $u_t$ is the per-period payoff. If all
sellers have cooperated (the buyer has received the agreed goods), then
the buyer resolves in finite time.
\end{theorem}

\begin{proof}
Let $\mathcal{B} = \sum_{i=1}^{n} 2\pay_i$ be the buyer's total
locked capital. Define the buyer's per-period payoff in each scenario:
\begin{itemize}
  \item \textbf{Resolve at time $t^*$}: Payoff at $t^*$ is
    $\sum_i \pi_{b,i} = \sum_i \pay_i$ (recovered capital).
    Payoff for $t < t^*$ is $0$ (capital locked, goods received).
    Total: $\delta^{t^*} \sum_i \pay_i$.
  \item \textbf{Never resolve}: Payoff at every $t$ is $0$ (capital
    locked forever, goods received but excess capital not recovered).
    Total: $0$.
\end{itemize}

The buyer prefers resolving at $t^*$ over never resolving iff:
\begin{equation}
  \delta^{t^*} \sum_i \pay_i > 0,
  \label{eq:liveness-condition}
\end{equation}
which holds for all finite $t^*$ since $\delta > 0$ and $\pay_i > 0$.

The buyer prefers resolving sooner (lower $t^*$) since
$\delta^{t_1} > \delta^{t_2}$ for $t_1 < t_2$ and $\delta < 1$.
Therefore the optimal strategy is $t^* = 0$ (resolve immediately).

A buyer who does not resolve at $t = 0$ forgoes
$\delta^0 \sum_i \pay_i = \sum_i \pay_i > 0$ in present value.
Since this exceeds the benefit of non-resolution (which is $0$ for
a rational agent with no taste for spite), indefinite non-resolution
is strictly dominated.
\end{proof}

\begin{remark}[Limitations]
\label{rem:liveness-limitations}
The result assumes: (i)~the buyer has received satisfactory performance
from all sellers, (ii)~the buyer has a positive discount rate
($\delta < 1$), and (iii)~the buyer does not derive positive utility
from harming sellers (no spite). Condition (iii) is the hardest to
guarantee. When it fails---an adversarial buyer who values
punishment above wealth---the mechanism cannot guarantee liveness
by economic incentives alone. This is the residual case addressed by
Layer~3 (legal deterrence via immutable on-chain evidence).
\end{remark}

\subsection{The Three-Layer Enforcement Architecture}

\begin{table}[ht]
\centering
\caption{Defense-in-depth: three enforcement layers.}
\label{tab:enforcement}
\begin{tabular}{llll}
\toprule
\textbf{Layer} & \textbf{Mechanism} & \textbf{Coverage} & \textbf{Failure Mode} \\
\midrule
1. Bonding  & Asymmetric collateral & $>$99\% of orders & Irrational actors \\
2. Coordination & Atomic resolution & Multi-seller failures & Uncoordinated sellers \\
3. Legal & Immutable evidence & Adversarial actors & Jurisdictional limits \\
\bottomrule
\end{tabular}
\end{table}

\paragraph{Layer 1: Economic.}
\Cref{thm:two-party-nash,thm:n-party-nash} establish that cooperation is
strictly dominant for all rational participants. This layer handles the
vast majority of transactions without any interaction beyond
commit-and-resolve.

\paragraph{Layer 2: Social.}
\Cref{prop:coordination-pressure} shows that atomic resolution creates
endogenous peer pressure. In multi-seller processes, honest sellers have
direct economic incentive to monitor and pressure underperforming
co-sellers. This layer requires no explicit communication protocol; the
incentive structure is sufficient.

\paragraph{Layer 3: Legal.}
For the residual case of truly irrational or adversarial actors, the
on-chain record provides tamper-proof evidence. Every commitment, every
bond deposit, every resolution (or lack thereof) is recorded with
block timestamps. This evidence can be presented in any off-chain legal
forum. Importantly, the protocol does \emph{not} attempt to perform
dispute resolution on-chain; it provides the \emph{evidence} for
off-chain systems that already exist.

% ════════════════════════════════════════════════════════════════════
\section{Implications for Organizational Theory}
\label{sec:post-firm}

\subsection{From Firms to Transaction-Scoped Institutions}

Coase's~\citeyearpar{coase1937} central insight is that firms exist
because market transaction costs---searching for partners, negotiating
terms, enforcing agreements---sometimes exceed the cost of hierarchical
management. Williamson~\citeyearpar{williamson1985} refined this into
Transaction Cost Economics (TCE), identifying three conditions under
which hierarchical governance dominates market exchange: high asset
specificity, high uncertainty, and high frequency.

Asymmetric bonding changes this calculus. Consider the three transaction
costs separately:

\paragraph{Search costs.}
Public on-chain signals (settlement history, accepted-token lists,
geolocation proofs) allow autonomous discovery through shared
coordination graphs. Search cost approaches zero for agents that can
read blockchain state.

\paragraph{Negotiation costs.}
EIP-712 off-chain commitment building allows parties to iterate on terms
(payment, agreement hash, token, topology) without gas costs. The
negotiation is bilateral and direct; counterparties can discover and
iterate through shared graph visibility rather than through a standing
intermediary's matching function.

\paragraph{Enforcement costs.}
This is the decisive variable. In the pre-Figaro world, enforcement
requires courts (slow, expensive, jurisdiction-dependent), arbitrators
(trusted third parties), or standing intermediaries that internalize
coordination and charge for it. \figaro{} prices enforcement at exactly
$2\times$ the transaction value---a fixed, predictable,
jurisdiction-independent cost requiring no third party.

\begin{proposition}[Coasean Threshold Shift]
\label{prop:coasean-shift}
For any transaction where $2\pay$ (the capital lockup cost) is less than
the firm's coordination overhead for the same transaction, the rational
organizational choice is market exchange via asymmetric bonding rather
than hierarchical integration.
\end{proposition}

The implication is structural: the boundary of the firm shifts
\emph{inward}, but the stronger implication is organizational rather
than merely contractual. When enforcement is priced directly in bonded
capital, the standing firm ceases to be the default container of
production. What emerges instead is a \emph{transaction-scoped
institution}: a temporary assembly of directly bonded counterparties
that forms around a process, coordinates for the duration of that
process, and dissolves at settlement. Functions previously internalized
for trust reasons---procurement, quality assurance, fulfillment
verification---become candidates for sovereign coordination between
independent value-adders. The ``restaurant'' becomes a coordination
label rather than a unitary firm: a cook, a kitchen operator, an
ingredient sourcer, and a courier each bond and settle independently
within the same process tree.

Each seller node in a process tree corresponds to what we term a
\emph{key component}---the irreducible asset or capability that
provides competitive advantage in a specific
market~\citep{daliana2013}. In the firm model, key components are
embedded inside organizational hierarchies: the cook's skill, the
kitchen's equipment, the courier's vehicle are all subsumed under the
firm's legal identity. In the \figaro{} model, each key component is
represented directly by a seller wallet address. The bonded process
tree does not merely flatten the firm---it makes each value-adding
capability individually addressable, individually compensated, and
individually accountable. The productive asset \emph{is} the identity.

Nothing in the mechanism requires that a node be a natural person or a
corporation. Any actor capable of controlling an address and bearing
capital risk---a human, a software agent, a machine operator, or a
treasury---can occupy a node in the process tree. The accountable unit
is the bonded counterparty, not the firm.

\subsection{Token Denomination as Coordination Signal}

In \figaro{}, the choice of denomination token is itself a coordination
signal. The protocol is token-agnostic: any ERC-20 can serve as
the bond and payment currency. This means:

\begin{itemize}
  \item A stablecoin signals legal-system alignment.
  \item A DAO governance token signals community membership.
  \item A commodity-backed token signals value anchoring.
  \item Accepting a specific token \emph{is} the seller's identity
    declaration---an accepted-token list is a brand.
\end{itemize}

Settlement velocity in a given token replaces traditional reputation
metrics (star ratings, review scores). The on-chain record of resolved
orders denominated in token $\tau$ is a public, tamper-proof signal of
the seller's reliability within the $\tau$ economy. No platform needs to
aggregate, curate, or attest this signal; it is a natural byproduct of
protocol usage.

\subsection{Cultural Hegemony and the Visibility of Coordination}

Gramsci~\citeyearpar{gramsci1971} observed that the most durable forms of
power are those embedded so deeply in common sense that they become
invisible. He called this \emph{cultural hegemony}: the dominant class
rules not only through coercion but by making its arrangements appear
natural, inevitable, and beyond question. The concept applies directly
to contemporary coordination infrastructure.

The platform economy illustrates Gramscian hegemony precisely. The
proposition that two strangers need a platform to transact safely feels
like common sense. It is not. It is a contingent arrangement that arose
because no economic primitive existed for safe direct exchange between
strangers. In Gramscian terms, the platform has achieved hegemonic
status: its necessity is no longer argued for---it is assumed.

Similarly, the proposition that production requires firms (the Coasean
assumption examined in \Cref{prop:coasean-shift}) and that cross-border
trade requires corridors controlled by state or alliance actors
(ports, railways, digital infrastructure) have both achieved the same
hegemonic status. These are presented as structural necessities rather
than contingent design choices.

Asymmetric bonding makes the coordination function \emph{visible} by
separating it from the entities that currently perform it. The platform's
coordination function (matching, enforcement, dispute resolution)
becomes a public smart contract. The firm's coordination function
(search, negotiation, enforcement) becomes three reducible costs. The
corridor's coordination function (settlement, standards, access
control) becomes a bilateral mechanism. In each case, the entity does
not need to be \emph{replaced}---it needs to be \emph{seen} as a design
choice rather than a fact of nature. Once the coordination function is
visible and separable, the architecture that makes the entity
structurally unnecessary follows.

This analysis explains a practical difficulty in communicating the
protocol's implications. Proposing to remove structures that most
people do not know are there is qualitatively different from proposing
a better version of the same structure. The challenge is not
engineering. It is making existing infrastructure legible as
architecture rather than gravity.

% ════════════════════════════════════════════════════════════════════
\section{Conclusion and Future Work}
\label{sec:conclusion}

We have presented \figaro{}, a smart-contract protocol that achieves
self-enforcing multi-party coordination through asymmetric bonding. The
mechanism is minimal (two external functions, no owner, no fee, no
escape hatches) and the game-theoretic properties are strong (cooperation
is strictly dominant at every position in an $N$-party process tree).

The implementation is formally verified through TLA\textsuperscript{+}
model checking (7~invariants, 2M+ states), property-based fuzzing, and
180~Foundry unit tests. The three-layer enforcement architecture---economic,
social, legal---provides defense-in-depth without on-chain governance or
dispute resolution.

Any associated token layer remains external to the kernel. In the
current runtime, seller-only settlement emission serves as a
coordination Schelling point; it is not required for bonding,
settlement, or participation.

\paragraph{Mechanism composition.}
The kernel's bonding equilibrium is a fixed point. The question---under
what conditions does an external mechanism composed with the kernel
preserve the Nash equilibrium?---is answered by the implementation's
convergence to the \emph{coordinator pattern}.

An external mechanism preserves the equilibrium if it satisfies four
constraints: (i)~read-only interaction with core state
(\texttt{staticcall} or event consumption), (ii)~no impersonation of
buyer or seller in core functions, (iii)~no modification to locked
bonds, and (iv)~no bypass of buyer-only resolution.

The reference implementation contains three proof cases. First, a Dutch
auction performs pure price discovery and role allocation with zero
token custody: the auction determines who may become the seller, but the
actual bond is still posted only through \texttt{commit}. Second, the
attestation coordinator appends schema-typed lifecycle and disclosure
events with role gating, but never touches locked capital or resolution
authority. Third, the schema registry and operator registry anchor
shared reference semantics and participant metadata as event-first
coordination layers; they are visible to agents and frontends, but
cannot alter the payoff matrix or bypass buyer-only resolution. All
compose with the kernel without weakening \Cref{thm:two-party-nash}.

Mechanism composition is therefore a layer above the kernel. The kernel
does not know or constrain what is composed around it. Declarative
institution assemblies specify which mechanisms compose for a given
workflow, with explicit risk classification (read-only, low-risk
coordinator, medium-risk extension, high-risk economic). The
coordinator pattern guarantees that read-only and event-only mechanisms
preserve the equilibrium by construction.

\paragraph{Broader implication.}
The Coasean threshold is not a conjecture---it is observable.
Contemporary coordination stacks often extract 15--30\% of transaction
value as coordination rent (Uber~25\%, DoorDash~30\%, Airbnb~14\%,
Amazon Marketplace~15\%).
Credit card networks layer 2--3\% interchange plus double-digit annual
interest on revolving balances. Financial markets intermediate
trillions in notional value against a fraction of that in real economic
exchange. Asymmetric bonding prices the cost of trust at $2\times$ the
transaction value in temporarily locked (not spent) capital, with zero
recurring extraction. For any economic activity where the coordination
rent exceeds the opportunity cost of locked capital, the standing
firm---as Coase conceived it---is no longer the uniquely efficient unit
of economic organization. The stronger consequence is the emergence of
transaction-scoped institutions: temporary assemblies of directly
bonded contributors that form around a process and dissolve at
settlement. The protocol does not abolish coordination; it changes the
unit of coordination from the standing firm to the bonded process.

\paragraph{Agent-native design.}
The protocol is designed for autonomous agent participation as a
first-class concern. EIP-712 typed commitments allow agents to build,
verify, and sign structured data using standard libraries. Event-sourced
state reconstruction means agents can derive the complete process graph
from on-chain logs without subgraph dependencies. The companion SDK
exposes a typed action proposer (\texttt{proposeActions}) that takes a
process state and an address, returning the set of available actions;
and a dual-mode execution queue (\texttt{ActionQueue}) that supports
both human-in-the-loop approval and fully autonomous submission. Nothing
in the mechanism distinguishes a human signer from an autonomous agent:
the bonding equilibrium holds for any rational actor capable of
controlling an address and bearing capital risk.

\paragraph{Execution environments.}
Future work includes verifying the kernel in alternative execution
environments, including the ongoing Cairo/Starknet port. The key
question is not only throughput, but which execution environments
preserve atomic resolution, buyer dominance, and evidence visibility
most cleanly without fragmenting a single process across multiple
settlement domains.

\section*{Acknowledgments}
I owe Vlad Zamfir a particular debt: his work in mechanism design and
coordination economics gave me the conceptual vocabulary to recognize
Buterin's Safe Remote Purchase contract not merely as an escrow pattern
but as a \emph{coordination mechanism}---one that could, in principle,
scale beyond two parties. Without that lens, the generalization to
asymmetric bonding and $N$-party process trees would not have
happened. I also want to acknowledge Vitalik Buterin for the Safe
Remote Purchase primitive itself, Fabrizio Romano Genovese for his
contributions to Figaro, and the Coalition of Automated Legal
Applications (COALA).


% ════════════════════════════════════════════════════════════════════
% References
% ════════════════════════════════════════════════════════════════════

\begin{thebibliography}{99}

\bibitem[Aragon(2017)]{aragon2017}
Cuende, L. and Izquierdo, J.
\newblock Aragon Network: A Decentralized Infrastructure for Value Exchange.
\newblock Aragon White Paper, 2017.
\newblock \url{https://github.com/aragon/whitepaper}

\bibitem[Asgaonkar and Krishnamachari(2019)]{asgaonkar2019}
Asgaonkar, A. and Krishnamachari, B.
\newblock Solving the Buyer and Seller's Dilemma: A Dual-Deposit Escrow Smart
  Contract for Provably Cheat-Proof Delivery and Payment for a Digital Good.
\newblock In \emph{Proc.\ IEEE International Conference on Blockchain and
  Cryptocurrency}, 2019.

\bibitem[BitHalo(2014)]{bithalo2014}
Dwyer, D.
\newblock BitHalo: Decentralized Contracts and Escrow Without Trust.
\newblock BitHalo White Paper, 2014.

\bibitem[Solidity Team(2016)]{buterin2016}
Solidity Team.
\newblock Safe Remote Purchase.
\newblock In \emph{Solidity Documentation --- Solidity by Example}, 2016.
\newblock \url{https://docs.soliditylang.org/en/latest/solidity-by-example.html#safe-remote-purchase}

\bibitem[Coase(1937)]{coase1937}
Coase, R.~H.
\newblock The Nature of the Firm.
\newblock \emph{Economica}, 4(16):386--405, 1937.

\bibitem[Daliana(2013)]{daliana2013}
Daliana, A.
\newblock \emph{ROKC: Leadership Built on the Return on Key Component}.
\newblock Chokti Inc., New York, 2013.

\bibitem[EIP-712(2017)]{eip712}
Bloemen, R., Logvinov, L., and Evans, J.
\newblock {EIP-712}: Ethereum typed structured data hashing and signing.
\newblock Ethereum Improvement Proposal, 2017.
\newblock \url{https://eips.ethereum.org/EIPS/eip-712}

\bibitem[Ghatak(1999)]{ghatak1999}
Ghatak, M.
\newblock Group lending, local information and peer selection.
\newblock \emph{Journal of Development Economics}, 60(1):27--50, 1999.

\bibitem[Grameen Bank(1999)]{grameen1999}
Yunus, M.
\newblock The Grameen Bank.
\newblock \emph{Scientific American}, 281(5):114--119, 1999.

\bibitem[Gramsci(1971)]{gramsci1971}
Gramsci, A.
\newblock \emph{Selections from the Prison Notebooks}.
\newblock International Publishers, New York, 1971.
\newblock Edited and translated by Q.~Hoare and G.~Nowell-Smith.

\bibitem[Hirshleifer(1983)]{hirshleifer1983}
Hirshleifer, J.
\newblock From weakest-link to best-shot: The voluntary provision of public
  goods.
\newblock \emph{Public Choice}, 41(3):371--386, 1983.

\bibitem[Kleros(2019)]{kleros2019}
Ast, F. and Lesaege, C.
\newblock Kleros: A Protocol for a Decentralized Justice System.
\newblock Kleros White Paper, 2019.
\newblock \url{https://kleros.io/yellowpaper.pdf}

\bibitem[Myerson(1981)]{myerson1981}
Myerson, R.~B.
\newblock Optimal auction design.
\newblock \emph{Mathematics of Operations Research}, 6(1):58--73, 1981.

\bibitem[Optimism(2021)]{optimism2021}
Optimism Foundation.
\newblock Optimistic Rollups.
\newblock Optimism Technical Documentation, 2021.
\newblock \url{https://community.optimism.io/docs/protocol/}

\bibitem[Poon and Dryja(2016)]{poon2016}
Poon, J. and Dryja, T.
\newblock The Bitcoin Lightning Network: Scalable Off-Chain Instant Payments.
\newblock Technical report, 2016.

\bibitem[Raiden(2017)]{raiden2017}
Raiden Network.
\newblock Raiden Network: Fast, Cheap, Scalable Token Transfers for Ethereum.
\newblock Raiden White Paper, 2017.
\newblock \url{https://raiden.network/}

\bibitem[UMA(2020)]{uma2020}
UMA Project.
\newblock The Optimistic Oracle.
\newblock UMA Technical Documentation, 2020.
\newblock \url{https://docs.uma.xyz/}

\bibitem[Williamson(1985)]{williamson1985}
Williamson, O.~E.
\newblock \emph{The Economic Institutions of Capitalism}.
\newblock Free Press, New York, 1985.

\end{thebibliography}

\end{document}
